\documentclass[11pt,a4paper]{article}

\usepackage[utf8]{inputenc}
\usepackage[T1]{fontenc}
\usepackage[spanish,es-noshorthands]{babel}
\usepackage{lmodern}
\usepackage[margin=2.3cm]{geometry}
\usepackage{xcolor}
\usepackage{listings}
\usepackage{booktabs}
\usepackage{tabularx}
\usepackage{enumitem}
\usepackage{array}
\usepackage[colorlinks=true,urlcolor=blue!60!black,linkcolor=black]{hyperref}

\definecolor{termbg}{gray}{0.95}
\lstdefinestyle{term}{
  basicstyle=\ttfamily\footnotesize,
  backgroundcolor=\color{termbg},
  frame=single, rulecolor=\color{gray!50},
  columns=fullflexible, keepspaces=true,
  breaklines=true, xleftmargin=4pt, xrightmargin=4pt,
  literate={á}{{\'a}}1 {é}{{\'e}}1 {í}{{\'i}}1 {ó}{{\'o}}1 {ú}{{\'u}}1
           {ñ}{{\~n}}1 {¡}{{!`}}1 {→}{{$\to$}}1
}
\lstset{style=term}
\newcommand{\cmd}[1]{\texttt{#1}}
\setlist{itemsep=2pt,topsep=4pt}

\title{\textbf{rvsim: Simulador de RISC-V (RV32I)}\\[2pt]
\large CS3051 Arquitectura de Computadores 2026-II --- Tarea 5}
\author{Ary Werner Aaron Rojas Durand}
\date{19 de noviembre de 2026}

\begin{document}
\maketitle

\section*{Enlaces de la entrega}
\begin{itemize}
  \item \textbf{Repositorio Git:} \url{https://github.com/MarianoUtec/ArquiHW5}
  \item \textbf{Carpeta de Google Drive} (video demostrativo + código fuente en ZIP):
        {No me alcanzó el tiempo a hacerlo, ya son las 11:50 pipipi}
\end{itemize}

\section{Resumen}
\textbf{rvsim} es un simulador del ISA base RISC-V de 32 bits (RV32I) con interfaz de terminal. Carga un programa
ya compilado, sea un binario crudo exportado de CPUlator o un ejecutable ELF, a partir de la dirección
\cmd{0x00000000}. Luego lo ejecuta instrucción por instrucción bajo control del usuario y permite inspeccionar en
cualquier momento el contador de programa (PC), los 32 registros y la memoria. No es un ensamblador: su entrada es
código máquina.

Cumple los cuatro requisitos obligatorios del enunciado e implementa tres de las cuatro ideas adicionales:
syscalls de SPIM, carga de ELF y desensamblador. También incluye breakpoints, un modo de traza y un mecanismo de
verificación automática que se usa en las pruebas.

\section{Lenguaje y organización del código}
El simulador está escrito en \textbf{C++17} estándar, sin dependencias externas, y compila con g++, clang o
MinGW mediante CMake o \cmd{make}. Se eligió un lenguaje de sistemas porque simular un ISA es, sobre todo,
manipular campos de bits y aritmética entera de 32 bits. Con \cmd{uint32\_t} el desborde es modular, igual que en
el hardware, e \cmd{int32\_t} se usa solo donde la instrucción interpreta los operandos con signo.

El código separa un \emph{núcleo}, que modela el procesador y no imprime nada, de la \emph{interfaz}:

\begin{center}
\begin{tabularx}{\textwidth}{@{}l l X@{}}
\toprule
Capa & Módulo (\cmd{src/}) & Responsabilidad \\
\midrule
Núcleo & \cmd{isa} & Decodificación: campos, inmediatos de los formatos I/S/B/U/J e identificación de la instrucción. \\
       & \cmd{cpu} & Estado arquitectural (PC, \cmd{x0}--\cmd{x31}) y ciclo fetch--decode--execute--write-back. \\
       & \cmd{memory} & Memoria plana little-endian de 32 bits, paginada. \\
       & \cmd{syscalls} & Servicios \cmd{ecall} al estilo SPIM (consola, \cmd{exit}, \cmd{sbrk}, tiempo). \\
       & \cmd{loader} & Carga de binarios crudos y de ELF32 (parser propio, sin \cmd{<elf.h>}). \\
Interfaz & \cmd{disasm} & Desensamblador con nombres ABI y pseudoinstrucciones. \\
       & \cmd{debugger} & Intérprete de comandos: ejecución, inspección, breakpoints, verificaciones. \\
       & \cmd{main} & Opciones de línea de comandos. \\
\bottomrule
\end{tabularx}
\end{center}

La CPU devuelve, por cada paso, un \cmd{Event} (normal, fin del programa, \cmd{ebreak}, error) y, si se le pide,
un \cmd{StepRecord} con lo que cambió: el registro destino con su valor anterior y nuevo, la escritura en memoria y
si hubo salto. La terminal decide cómo mostrarlo. Las syscalls se conectan a través de la interfaz
\cmd{SyscallHandler}, así que el núcleo podría reutilizarse con otra interfaz sin modificarse.

\subsection{Ciclo de ejecución}
\cmd{Cpu::step()} ejecuta una instrucción en cuatro etapas:
\begin{enumerate}
  \item \textbf{Fetch:} lee los 4 bytes en \cmd{pc} y los arma en orden little-endian. Si el PC no está alineado a
        4, se reporta un error.
  \item \textbf{Decode:} extrae \cmd{opcode}, \cmd{rd}, \cmd{funct3}, \cmd{rs1}, \cmd{rs2} y \cmd{funct7}, arma el
        inmediato según el formato con extensión de signo e identifica la instrucción exacta. Las combinaciones que
        no son RV32I se marcan como \emph{ilegales} con un motivo: \cmd{funct7} inválido, \cmd{slli} con
        \cmd{imm[11:5]}~$\neq 0$, instrucciones de la extensión M, instrucciones comprimidas o la palabra
        \cmd{0x00000000}.
  \item \textbf{Execute:} un \cmd{switch} sobre la instrucción calcula el resultado, la dirección efectiva de
        loads y stores o la condición del branch.
  \item \textbf{Write-back y PC:} escribe \cmd{rd} (nunca \cmd{x0}) y actualiza \cmd{pc} a $pc+4$ o al destino del
        salto. En \cmd{jalr} se limpia el bit 0 del destino, y un destino no alineado a 4 produce el error
        \emph{instruction-address-misaligned}.
\end{enumerate}

\subsection{Memoria}
El espacio de direcciones es de $2^{32}$ bytes, direccionable por byte y little-endian. Como reservar 4\,GiB no es
viable, la memoria se guarda como un \cmd{unordered\_map} de páginas de 4\,KiB que se crean la primera vez que se
escriben. Leer una dirección que nunca se escribió devuelve 0 sin crear la página. Los accesos de 16 y 32 bits se
arman byte a byte, de modo que el orden de bytes es explícito y funcionan aunque crucen el borde de una página.

\section{Características}
\subsection*{Requisitos obligatorios}
\begin{itemize}
  \item \textbf{Las 37 instrucciones del Anexo A.}
        Loads: \cmd{lb lh lw lbu lhu}.
        OP-IMM: \cmd{addi slli slti sltiu xori srli srai ori andi}.
        Upper: \cmd{lui auipc}.
        Stores: \cmd{sb sh sw}.
        OP: \cmd{add sub sll slt sltu xor srl sra or and}.
        Branches: \cmd{beq bne blt bge bltu bgeu}.
        Saltos: \cmd{jal jalr}.
        Además se reconocen \cmd{ecall}, \cmd{ebreak} y \cmd{fence} (como \emph{nop}).
  \item \textbf{Memoria little-endian} de 32 bits.
  \item \textbf{Carga de binarios crudos} en \cmd{0x00000000} (la dirección es configurable con \cmd{--base}).
  \item \textbf{Ejecución paso a paso} con \cmd{step [n]}; \emph{Enter} vacío repite el último paso.
  \item \textbf{Inspección} de PC (\cmd{pc}), registros (\cmd{regs}) y memoria (\cmd{mem}, \cmd{words}, \cmd{dump}).
\end{itemize}

\subsection*{Ideas adicionales implementadas}
\begin{itemize}
  \item \textbf{Syscalls de SPIM} vía \cmd{ecall}, con el servicio en \cmd{a7} como en CPUlator:
        1 print\_int, 4 print\_string, 5 read\_int, 8 read\_string, 9 sbrk, 10 exit, 11 print\_char,
        12 read\_char, 17 exit2, 30 time, 32 sleep, 34 print\_hex, 35 print\_bin y 36 print\_uint. También se acepta
        el 93 (exit de Linux), que usan los programas en C.
  \item \textbf{Carga de ELF32}: segmentos \cmd{PT\_LOAD} (con relleno en cero de \cmd{.bss}), punto de entrada
        \cmd{e\_entry} y tabla de símbolos. Los símbolos se usan en breakpoints (\cmd{break main}), en expresiones
        (\cmd{words primos 10}) y en el desensamblado (\cmd{<criba+12>}).
  \item \textbf{Desensamblador}: cada \cmd{step} muestra la instrucción en ejecución, y \cmd{dis} lista cualquier
        rango, con nombres ABI y la pseudoinstrucción equivalente (\cmd{li}, \cmd{mv}, \cmd{j}, \cmd{call},
        \cmd{ret}, \cmd{beqz}, \ldots).
\end{itemize}
No se implementó el soporte para GDB.

\subsection*{Otras funciones}
\begin{itemize}
  \item Breakpoints por dirección o símbolo; \cmd{ebreak} pausa la ejecución y Ctrl+C detiene un \cmd{run} largo.
  \item \cmd{trace on}: durante \cmd{run} se muestra cada instrucción con sus efectos.
  \item Errores explícitos: instrucción ilegal (con su motivo), PC o salto desalineado, syscall desconocida y
        entrada agotada. La CPU queda detenida y \cmd{reset} recarga el programa desde el archivo.
  \item Detección del fin típico de los programas de CPUlator (\cmd{end: j end}, un salto a sí mismo).
  \item Modos sin terminal: \cmd{--run}, scripts (\cmd{-x}) y comandos en línea (\cmd{-c}), con verificaciones
        \cmd{expect} para pruebas automáticas.
\end{itemize}

\section{Uso}
\begin{lstlisting}
cmake -B build && cmake --build build      # o: make
./build/rvsim programas/riscvtest.bin       # terminal interactiva
./build/rvsim -r programas/quicksort.bin    # ejecutar de corrido
\end{lstlisting}

\begin{center}
\begin{tabularx}{\textwidth}{@{}l X@{}}
\toprule
Comando & Descripción \\
\midrule
\cmd{step [n]} & Ejecuta $n$ instrucciones (1 por defecto) mostrando sus efectos. \\
\cmd{run [max]} & Ejecuta hasta el fin del programa, un breakpoint, \cmd{ebreak}, un error o \cmd{max} instrucciones. \\
\cmd{break <dir|símbolo>}, \cmd{delete [dir]} & Pone o quita breakpoints. \\
\cmd{pc [valor]} & Muestra o cambia el PC. \\
\cmd{regs [r \ldots]} & Los 32 registros (marca con \cmd{*} los modificados) o solo los indicados. \\
\cmd{mem <ini> [fin]} & Bytes de memoria en el rango, en el formato del ejemplo del enunciado. \\
\cmd{words}, \cmd{dump}, \cmd{str} & Palabras de 32 bits, volcado hexadecimal + ASCII o una cadena. \\
\cmd{dis [dir] [n]} & Desensambla $n$ instrucciones. \\
\cmd{set}, \cmd{setmem} & Modifican un registro o la memoria. \\
\cmd{info}, \cmd{reset}, \cmd{load}, \cmd{help}, \cmd{exit} & Estado general, recarga, otro programa, ayuda, salir. \\
\bottomrule
\end{tabularx}
\end{center}

\noindent Las direcciones y los valores aceptan \cmd{4096}, \cmd{0x1000}, \cmd{-1}, un registro (\cmd{sp}), un
símbolo (\cmd{main}) o una suma o resta simple (\cmd{sp+8}). Ejemplo de sesión:

\begin{lstlisting}
$ ./build/rvsim programas/riscvtest.bin
"programas/riscvtest.bin" cargado a memoria. Escribe 'help' para ver los comandos.
> step
   0x00000000:  00500113  addi   sp, zero, 5          # li sp, 5
                          sp: 0x00000000 -> 0x00000005  (5)
> run
Fin del programa: salto a sí mismo en 0x00000050 (bucle final del programa).
Ejecutadas 18 instrucciones. PC = 0x00000050
> mem 0x64 0x67
Memoria (0x64-0x67): 19 00 00 00
\end{lstlisting}

\section{Verificación}
El simulador se probó con los tres programas que recomienda el enunciado y con cuatro programas propios. Cada uno
tiene un script en \cmd{tests/} que verifica el resultado con \cmd{expect}, y todos pasan
(\cmd{make test} o \cmd{ctest}).

\begin{center}
\begin{tabularx}{\textwidth}{@{}l X l@{}}
\toprule
Programa & Qué prueba & Resultado \\
\midrule
\cmd{riscvtest.bin} & add, sub, and, or, slt, addi, lw, sw, beq, jal & \cmd{mem[0x64] = 25} \\
\cmd{quicksort.bin} & Recursión, pila, loads y stores sobre un arreglo & \{1,2,3,4,6,8,9\} en \cmd{0x1000} \\
\cmd{arbol.bin} & \cmd{is\_symmetric} sobre \{6,4,4,1,2,2,1\} & \cmd{a0 = 1} \\
\cmd{rv32i\_suite.bin} & Las 37 instrucciones, 46 casos (signo, overflow, shifts $\geq 32$, offsets negativos) & 46/46 OK \\
\cmd{fibonacci.bin} & Syscalls de lectura y escritura; serie guardada en memoria & \cmd{fib(10) = 55} \\
\cmd{primos.elf} & Programa en C cargado como ELF, breakpoint por símbolo & primos $\leq N$ \\
\cmd{invalida.bin} & \cmd{mul} (extensión M) dentro de un programa & error explícito en \cmd{0xC} \\
\bottomrule
\end{tabularx}
\end{center}

Además, durante el desarrollo se hicieron dos pruebas diferenciales:
\begin{itemize}
  \item \textbf{Decodificador:} 3000 instrucciones RV32I aleatorias, de todos los formatos y con inmediatos y
        registros al azar, ensambladas con LLVM. Su desensamblado coincidió en el 100\,\% de los casos con el de
        \cmd{llvm-objdump}.
  \item \textbf{Ejecución:} 150 programas aleatorios de 60 instrucciones (aritméticas, lógicas, shifts, loads y
        stores) se ejecutaron en el simulador y en un modelo de referencia independiente escrito en Python. Los
        32 registros finales coincidieron en todos los casos.
\end{itemize}

\section{Limitaciones}
Solo se implementa el ISA base RV32I: las extensiones M, A, F, D y C y las instrucciones CSR se reportan como
ilegales. No hay modos de privilegio, excepciones con \emph{trap handler} ni soporte para GDB. Los accesos a
memoria no alineados se permiten y se resuelven byte a byte, una de las opciones válidas según la especificación.

\end{document}
